\documentclass[11pt,a4paper]{article}

% ============================================================
% PACKAGES
% ============================================================
\usepackage[utf8]{inputenc}
\usepackage[T1]{fontenc}
\usepackage[english]{babel}
\usepackage{amsmath,amssymb,amsthm,mathtools}
\usepackage{mathrsfs}
\usepackage{tikz-cd}
\usepackage{enumitem}
\usepackage{hyperref}
\usepackage{geometry}
\usepackage{microtype}
\usepackage{booktabs}
\usepackage{csquotes}

\geometry{margin=1in}

% ============================================================
% THEOREM ENVIRONMENTS (lighter — this is a position paper)
% ============================================================
\theoremstyle{definition}
\newtheorem{definition}{Definition}[section]
\newtheorem{principle}[definition]{Principle}
\newtheorem{example}[definition]{Example}
\newtheorem{remark}[definition]{Remark}
\newtheorem{objection}[definition]{Objection}
\newtheorem{reply}[definition]{Reply}

\theoremstyle{plain}
\newtheorem{proposition}[definition]{Proposition}
\newtheorem{thesis}[definition]{Thesis}

% ============================================================
% MACROS (inherited, light)
% ============================================================
\newcommand{\Ccal}{\mathcal{C}}
\newcommand{\Xcal}{\mathcal{X}}
\newcommand{\Rcal}{\mathcal{R}}
\newcommand{\Tcal}{\mathcal{T}}
\newcommand{\Dcal}{\mathcal{D}}
\newcommand{\Ecal}{\mathcal{E}}
\newcommand{\Vcal}{\mathcal{V}}
\newcommand{\Hcal}{\mathcal{H}}
\newcommand{\Acal}{\mathcal{A}}
\newcommand{\Norm}[1]{\left\lVert #1 \right\rVert}
\newcommand{\R}{\mathbb{R}}
\newcommand{\N}{\mathbb{N}}
\newcommand{\Z}{\mathbb{Z}}
\newcommand{\op}[1]{\operatorname{#1}}
\newcommand{\Phimap}{\Phi}

% ============================================================
% TITLE
% ============================================================
\title{%
  \textbf{Structural Auditability in Mathematical Construction:}\\[4pt]
  \large A Methodological Position
}
\author{}
\date{}

\begin{document}
\maketitle

\begin{abstract}
\noindent
We argue that a central methodological question in mathematical
practice is not only \emph{whether} a construction is correct, but
\emph{whether it is auditable}---that is, whether each step of the
construction contains within itself the means to anticipate and
localise objections. Building on two companion papers on a
multirepresentational constraining methodology, we formalise
\emph{structural wear} as the mechanism by which a construction is
successively tested against independent representations, and we
introduce the \emph{certification trace} as the record of this
process. We develop a multi-level taxonomy of methodological
outcomes, distinguishing proof, formalisation, extension,
characterisation, numerical evidence with identified limits, and
structural obstruction. We argue that this taxonomy replaces the
binary success/failure framework with a finer-grained account of
what a construction delivers. We close with a set of objections and
replies, and with a discussion of the relationship between
auditability and external peer review.
\end{abstract}

\tableofcontents

% ------------------------------------------------------------
\section{From Proof to Architecture}
\label{sec:intro}
% ------------------------------------------------------------

\subsection{The Standard Picture}

The standard picture of mathematical work is broadly deductive:
a \emph{definition} is followed by a \emph{lemma}, which is
followed by a \emph{theorem}, which is followed by a
\emph{corollary}. Correctness is measured by the validity of the
deductive chain, and the chain is audited by peer review. If the
chain contains a gap, the reviewer identifies it and the author
either closes it or retracts the claim.

This picture is extraordinarily effective. It has produced
mathematics of immense depth, and it is the operational standard
against which formal work is judged.

\subsection{A Complementary Question}

The picture is, however, incomplete as a description of
\emph{how mathematical knowledge is constructed}. Before the
deductive chain exists, there is a phase of construction in which
representations are chosen, transformations are applied, and the
shape of the eventual theorem is progressively determined. The
question that governs this phase is not

\begin{quote}
\emph{Is this step valid?}
\end{quote}

but

\begin{quote}
\emph{If this step is wrong, how will I know?}
\end{quote}

This is the question of \emph{auditability}. It is not a
replacement for deductive rigour; it is the design principle that
governs how a construction should be organised so that deductive
rigour can eventually be applied.

\subsection{Thesis}

\begin{thesis}[Auditability]\label{th:audit}
A mathematical construction should be designed so that each step
contains the means to audit the previous step. The strength of a
construction is measured not only by the correctness of its final
claim, but by the redundancy, coherence, and auditability of the
chain that produces that claim.
\end{thesis}

\subsection{Plan}

Section~\ref{sec:audit-classical} analyses the audit relationship
in classical mathematics. Section~\ref{sec:wear} introduces
\emph{structural wear} as the mechanism of auditability.
Section~\ref{sec:redundancy} develops the notions of redundancy and
coherence. Section~\ref{sec:taxonomy} presents the multi-level
taxonomy of methodological outcomes. Section~\ref{sec:trace}
introduces the certification trace.
Section~\ref{sec:cycle} describes the self-auditing cycle.
Section~\ref{sec:objections} treats objections and replies.
Section~\ref{sec:relationship} relates the position to two
companion papers \cite{paperI,paperII}.

% ------------------------------------------------------------
\section{The Audit Relationship}
\label{sec:audit-classical}
% ------------------------------------------------------------

\subsection{The Classical Model}

In the classical model, the relationship between author and
auditor is:
\[
  \text{author} \;\longrightarrow\; \text{construction}
  \;\longrightarrow\; \text{auditor} \;\longrightarrow\;
  \begin{cases} \text{accept} \\ \text{reject} \end{cases}
\]
The auditor's role is adversarial: to find gaps, hidden
hypotheses, unjustified steps, and counterexamples. The author's
role is defensive: to present a construction that resists the
auditor's attacks.

\subsection{The Cost of the Classical Model}

The classical model has two costs.

First, it is \emph{post hoc}: the auditor sees only the final
construction, not the process. Hidden assumptions that were
resolved internally are invisible; unresolved ones appear as
gaps.

Second, it is \emph{asymmetric}: the author knows the construction
better than the auditor, but has no obligation to expose the
structure of the argument. A well-written paper can hide a weak
foundation behind elegant presentation.

\subsection{What Auditability Adds}

Auditability adds a design constraint on the author:
\[
  \boxed{
    \text{Construct as if the auditor were present at every step.}
  }
\]
The author is not required to eliminate the auditor. The author is
required to \emph{anticipate} the auditor's questions and to
provide, at each step, the means to answer them.

\begin{principle}[Anticipated audit]\label{pr:anticipated}
Each step of a construction should be presented together with the
questions a competent auditor would ask about it and the answers
available at that step.
\end{principle}

% ------------------------------------------------------------
\section{Structural Wear}
\label{sec:wear}
% ------------------------------------------------------------

\subsection{The Mechanism}

A construction is not made stronger by adding more of the same
representation. It is made stronger by subjecting it to
\emph{structurally diverse} representations, each of which tests a
different aspect of the construction.

\begin{definition}[Structural wear]\label{def:wear}
Given a construction $X$ and a family of representations
$\Rcal = \{R_i\}_{i \in I}$, the \emph{structural wear} of $X$
under $\Rcal$ is the process
\[
  X
  \;\xrightarrow{R_1}\;
  X_1
  \;\xrightarrow{R_2}\;
  X_2
  \;\longrightarrow\;
  \cdots
  \;\xrightarrow{R_n}\;
  X_n
\]
together with the requirement that
\[
  X_k \simeq X
  \qquad \text{(structurally, for each } k \text{)}.
\]
\end{definition}

\subsection{What Wear Eliminates}

Wear does not destroy the object. It eliminates the degrees of
freedom that are not compatible with the represented constraints.
In the language of the companion papers \cite{paperI,paperII}:
\[
  \Ccal_0 \supseteq \Ccal_1 \supseteq \cdots \supseteq \Ccal_n,
\]
and the surviving structure is
\[
  \Ccal_\infty = \bigcap_i \Ccal_i .
\]

\subsection{Why the Metaphor of Wear is Appropriate}

The metaphor is appropriate because each representation is
\emph{hostile} to the construction. It does not confirm; it tests.
The construction must survive the test, and survival is measured
by structural identity:
\[
  R_i(X) \neq \varnothing
  \quad\text{and}\quad
  R_i^{-1}(R_i(X)) \simeq X.
\]

\begin{principle}[Hostile representation]\label{pr:hostile}
Each representation should be applied as if its purpose were to
destroy the construction. The construction is retained only if it
survives.
\end{principle}

% ------------------------------------------------------------
\section{Redundancy and Coherence}
\label{sec:redundancy}
% ------------------------------------------------------------

\subsection{Three Types of Redundancy}

Redundancy in mathematical construction has three levels.

\begin{definition}[Weak redundancy]\label{def:weak-red}
Representations $R_1, R_2, \ldots, R_k$ are \emph{weakly
redundant} if they transmit essentially the same information
expressed in different languages.
\end{definition}

\begin{definition}[Structural redundancy]\label{def:strong-red}
Representations are \emph{structurally redundant} if they impose
different and jointly compatible constraints on the same object.
\end{definition}

\begin{definition}[Cross redundancy]\label{def:cross-red}
Representations are \emph{cross redundant} if there exist mutual
implications
\[
  R_i \Rightarrow R_j
  \quad\text{and}\quad
  R_j \Rightarrow R_i
\]
in the sense that a constraint from one representation can be
derived from a constraint of the other.
\end{definition}

\subsection{The Strength Ordering}

The strength of a construction increases with the type of
redundancy employed:
\[
  \text{weak} \;<\; \text{structural} \;<\; \text{cross}.
\]

\begin{proposition}[Cross redundancy as auditability]\label{pr:cross-audit}
Cross redundancy is the strongest form of auditability available
within a representation family, because it allows each
representation to check the other without external intervention.
\end{proposition}

\subsection{Coherence}

Coherence is the requirement that the representations of a single
object are mutually consistent. In the language of the companion
paper \cite{paperII}:
\[
  \Rcal_{\rm coh} \subset \prod_i \Rcal_i .
\]
The admissible space of representations is a subspace, not a
product: some combinations of individual representations are
compatible, others are not.

\begin{principle}[Coherence as constraint]\label{pr:coherence}
The representations of a construction should be chosen so that
their compatibility relations constrain the object beyond what any
individual representation determines.
\end{principle}

% ------------------------------------------------------------
\section{The Multi-Level Output Taxonomy}
\label{sec:taxonomy}
% ------------------------------------------------------------

\subsection{The Binary Framework and its Limits}

The classical framework classifies a construction as either
\emph{correct} or \emph{incorrect}. This binary framework is
adequate for finished theorems but inadequate for constructions in
progress.

\begin{remark}
A construction that does not yield a proof may nonetheless yield
substantial mathematical knowledge: a rigorous characterisation of
an obstruction, an identified bottleneck, a set of numerical
evidence with controlled error, or an extension of an existing
theory.
\end{remark}

\subsection{The Six Levels}

\begin{definition}[Output taxonomy]\label{def:taxonomy}
The outcome of a construction falls into one of six levels:
\begin{enumerate}[label=\textbf{Level \Roman*.}]
  \item \textbf{Complete proof:} the construction closes with a
        theorem and its proof.
  \item \textbf{Rigorous formalisation:} the construction is
        expressed in a formal system with all definitions, lemmas,
        and derivations explicit.
  \item \textbf{Extension of existing theory:} the construction
        produces a generalisation or strengthening of a known
        result.
  \item \textbf{Structural characterisation:} the construction
        determines the object up to an explicitly declared
        equivalence.
  \item \textbf{Numerical evidence with identified limits:} the
        construction produces controlled numerical evidence
        together with an explicit statement of where the evidence
        stops and what would be required to continue.
  \item \textbf{Structural obstruction:} the construction
        identifies an obstruction that prevents further progress,
        and localises the obstruction at a specific layer,
        criterion, or hypothesis.
\end{enumerate}
\end{definition}

\subsection{Why Six Levels}

The six levels are not a hierarchy of decreasing value. They are a
classification of \emph{what kind of knowledge} the construction
produces.

\begin{proposition}[Non-subordination]\label{pr:non-sub}
Level~\RomanVI{} (structural obstruction) is not inferior to
Level~\RomanI{} (complete proof). A rigorous identification of an
obstruction is often more valuable than a proof of a special case,
because it explains \emph{why} the general case resists proof.
\end{proposition}

\begin{remark}
This is the formal content of the informal observation that
``knowing why a problem is hard'' is itself a mathematical
achievement. The taxonomy makes this observation precise.
\end{remark}

\subsection{Examples}

\begin{example}[Level~\RomanI]
The Dirichlet Laplacian on $[0,1]$ closes with an explicit
spectral theorem. Complete proof.
\end{example}

\begin{example}[Level~\RomanIV]
The classification of surfaces up to homeomorphism is a structural
characterisation: the object is determined up to the declared
equivalence.
\end{example}

\begin{example}[Level~\RomanVI]
The failure of injectivity of the inverse spectral map (Gordon--Webb--Wolpert)
is a structural obstruction: it localises the failure of
identifiability at the spectral layer.
\end{example}

% ------------------------------------------------------------
\section{The Certification Trace}
\label{sec:trace}
% ------------------------------------------------------------

\subsection{Motivation}

A finished construction carries with it a history: definitions
were chosen, hypotheses were declared, representations were
applied, restrictions were imposed, evidence was gathered,
verifications were performed, and corrections were made. In the
classical framework, this history is largely invisible.

\begin{definition}[Certification trace]\label{def:trace}
The \emph{certification trace} of a construction $X$ is the tuple
\[
  \Tcal(X)
  =
  (D, H, C, R, E, V, A),
\]
where
\begin{itemize}
  \item $D$ = the set of definitions used;
  \item $H$ = the set of hypotheses declared;
  \item $C$ = the set of constructions performed;
  \item $R$ = the set of representations applied;
  \item $E$ = the set of estimates and evidence gathered;
  \item $V$ = the set of cross-verifications performed;
  \item $A$ = the set of audits and corrections recorded.
\end{itemize}
\end{definition}

\subsection{What the Trace Records}

The trace is not a proof. It is a record of the process by which
the proof was produced. It is intended to answer, for any step of
the final construction, the questions:
\begin{itemize}
  \item Where did this definition come from?
  \item What hypothesis does this step rely on?
  \item Which representation is being applied?
  \item What evidence supports this step?
  \item Which other representation confirms it?
  \item Has this step been audited?
\end{itemize}

\begin{principle}[Trace completeness]\label{pr:trace}
A construction is \emph{trace-complete} if every step of the final
proof is covered by at least one entry in the certification trace.
\end{principle}

\subsection{Trace as Math Object}

\begin{remark}
The certification trace is itself a mathematical object: it is a
finite structured tuple that can be manipulated, compared, and
formalised. In particular, it can be encoded in a proof assistant,
where the entries become proof obligations.
\end{remark}

% ------------------------------------------------------------
\section{The Self-Auditing Cycle}
\label{sec:cycle}
% ------------------------------------------------------------

\subsection{The Cycle}

The self-auditing cycle is:
\[
  \boxed{
    \text{construct}
    \to
    \text{represent}
    \to
    \text{constrain}
    \to
    \text{audit}
    \to
    \text{correct}
    \to
    \text{reconstruct}.
  }
\]

Each arrow is a distinct operation:
\begin{itemize}
  \item \textbf{construct}: produce a candidate model;
  \item \textbf{represent}: apply a layer representation;
  \item \textbf{constrain}: intersect with the admissible model space;
  \item \textbf{audit}: test for cross-coherence;
  \item \textbf{correct}: localise and repair inconsistencies;
  \item \textbf{reconstruct}: return to the original object via
        inverse representation.
\end{itemize}

\subsection{Correction as Information}

\begin{principle}[Correction as information]\label{pr:correction}
When a contradiction between representations is found, it is
recorded as a structural feature of the construction, not hidden.
The contradiction localises the missing information.
\end{principle}

\begin{proposition}[Diagnostic value of contradiction]\label{pr:diagnostic}
If representations $R_i$ and $R_j$ fail to be cross-coherent on a
construction $X$, the failure identifies a specific structural
feature of $X$ that one of the representations does not capture.
\end{proposition}

\subsection{Relationship to External Audit}

The self-auditing cycle does not eliminate the external auditor.
It changes the state of the construction at the time of external
audit:

\begin{quote}
\emph{Classical:} the auditor sees a finished construction and
must find its flaws.\\[2mm]
\emph{Self-audited:} the auditor sees a construction whose
anticipated flaws have been addressed, and must find the
remaining ones.
\end{quote}

\begin{remark}
The self-auditing cycle is a \emph{reduction} of the external
auditor's workload, not a substitution for it.
\end{remark}

% ------------------------------------------------------------
\section{Objections and Replies}
\label{sec:objections}
% ------------------------------------------------------------

\subsection{Objection 1: This is Just Careful Mathematics}

\begin{objection}
Everything described here is what good mathematicians already do:
choose diverse representations, check for consistency, identify
where proofs fail.
\end{objection}

\begin{reply}
The observation is correct. The contribution of this position paper
is not the invention of a new practice, but the \emph{explicit
formalisation} of a practice that is currently implicit. Making the
practice explicit allows:
\begin{itemize}
  \item (a) the design of a \emph{certification trace} that can be
        mechanically checked;
  \item (b) the definition of \emph{multi-level outputs} that
        classify partial results;
  \item (c) the \emph{quantification} of redundancy and coherence.
\end{itemize}
Implicit practice cannot be formalised; explicit practice can.
\end{reply}

\subsection{Objection 2: This Does Not Prove Anything}

\begin{objection}
The methodology produces evidence, characterisations, and
obstructions. None of these is a proof. Therefore the methodology
does not contribute to mathematical knowledge.
\end{objection}

\begin{reply}
The premise is that only proofs contribute to mathematical
knowledge. The position paper argues the contrary: the \emph{levels
taxonomy} of \S\ref{sec:taxonomy} classifies six distinct types of
contribution, of which complete proof is only one. A rigorous
characterisation of an obstruction (Level~\RomanVI) is a
mathematical contribution even when no proof exists.
\end{reply}

\subsection{Objection 3: Auditability is Not Measurable}

\begin{objection}
``Auditability'' is a vague term. Without a quantitative measure,
the principle cannot be applied.
\end{objection}

\begin{reply}
Auditability can be made quantitative via the \emph{certification
trace} $\Tcal(X)$. A construction is trace-complete (Principle~\ref{pr:trace})
if every step of the final proof is covered by an entry in the
trace. The trace is finite and structured, and can be counted,
compared, and formalised. Its completeness is a checkable property.
\end{reply}

\subsection{Objection 4: This Does Not Apply to All Mathematics}

\begin{objection}
Some mathematics does not admit multiple representations. The
methodology is inapplicable in such cases.
\end{objection}

\begin{reply}
The methodology is a \emph{design principle}, not a universal
requirement. It applies when multiple representations are available
and cross-coherent. In contexts where only one representation is
available, the methodology reduces to the classical framework. The
position paper does not claim universality.
\end{reply}

\subsection{Objection 5: Why Bother Formalising Implicit Practice?}

\begin{objection}
If the practice is already effective, why formalise it?
\end{objection}

\begin{reply}
Because formalisation enables three things that implicit practice
cannot achieve:
\begin{itemize}
  \item \emph{Teaching}: implicit practice can only be transmitted
        by apprenticeship; explicit practice can be taught.
  \item \emph{Mechanical verification}: implicit practice cannot be
        encoded in a proof assistant; explicit practice can.
  \item \emph{Extension}: implicit practice cannot be systematically
        extended; explicit practice can be generalised to new
        domains.
\end{itemize}
\end{reply}

% ------------------------------------------------------------
\section{Relationship to Companion Papers}
\label{sec:relationship}
% ------------------------------------------------------------

\subsection{The Triad}

The three papers form a triad:

\begin{center}
\begin{tabular}{c|c|c}
  \toprule
  Paper & Question & Answer \\
  \midrule
  \cite{paperI} & How does the method work?
    & Constraining, spectral emergence, discrepancy classification \\
  \cite{paperII} & When does it fail?
    & Separation, three-type rupture, coherence subspace \\
  This paper & What does success mean?
    & Auditability, wear, taxonomy, trace \\
  \bottomrule
\end{tabular}
\end{center}

\subsection{Independence of the Papers}

Each paper is independent of the others:
\begin{itemize}
  \item Paper~\cite{paperI} can be read as a methodology of
        construction with concrete toys.
  \item Paper~\cite{paperII} can be read as a meta-theory of the
        limits of constraining.
  \item This paper can be read as an epistemological position on
        what mathematical construction should be.
\end{itemize}
Together they form a coherent programme.

\subsection{The Shared Vocabulary}

The three papers share:
\begin{itemize}
  \item the master object $\mathfrak O$;
  \item the layer set $\Lambda$;
  \item the transition morphisms $\phi_{\lambda\mu}$;
  \item the validation matrix $\mathbf V$;
  \item the epistemic states $\{P, N, O, F\}$.
\end{itemize}
This shared vocabulary is the technical backbone of the programme.

% ------------------------------------------------------------
\section{Conclusion}
\label{sec:conclusion}
% ------------------------------------------------------------

The position defended in this paper is:

\begin{thesis}[Restated]
Mathematical construction should be designed so that each step
contains the means to audit the previous step. Success is not
binary: it is a multi-level outcome classified by the taxonomy of
\S\ref{sec:taxonomy}.
\end{thesis}

The mechanism by which auditability is achieved is \emph{structural
wear}: successive applications of independent, cross-coherent
representations, each of which tests the construction against a
different aspect of its structure.

The record of the process is the \emph{certification trace}, a
finite structured object that can be manipulated, compared, and
formalised.

The self-auditing cycle is the operational realisation of the
principle: it does not replace external audit, but reduces its
workload by anticipating its questions.

\medskip

The methodological slogan, in one line:

\begin{quote}
\emph{Build so that the auditor is present at every step.}
\end{quote}

Or, more precisely:

\begin{quote}
\emph{A mathematical construction is well-designed when each of its
steps contains the means to anticipate and answer the questions a
competent auditor would ask about it.}
\end{quote}

% ------------------------------------------------------------
\appendix
% ------------------------------------------------------------

\section{Summary of Principles}
\label{app:principles}

\begin{center}
\begin{tabular}{c l}
  \toprule
  Principle & Statement \\
  \midrule
  \ref{pr:anticipated} & Anticipated audit \\
  \ref{pr:hostile} & Hostile representation \\
  \ref{pr:coherence} & Coherence as constraint \\
  \ref{pr:trace} & Trace completeness \\
  \ref{pr:correction} & Correction as information \\
  \bottomrule
\end{tabular}
\end{center}

\section{Bibliography}
\label{app:biblio}

\begin{thebibliography}{99}

\bibitem{paperI}
[Companion paper I],
\emph{Structural Constraining and Spectral Emergence: A
Multirepresentational Methodology for Nonlinear Evolution
Equations}, 2024.

\bibitem{paperII}
[Companion paper II],
\emph{On the Limits of Structural Constraining: Separation,
Coherence, and the Emergence of New Representations}, 2024.

\bibitem{framework}
[Framework paper],
\emph{A Multirepresentational Construction--Reconstruction
Methodology for Mathematical Structures}, Parts I--V, 2024.

\bibitem{polya}
G.~Polya,
\emph{How to Solve It}, Princeton University Press, 1945.

\bibitem{lakatos}
I.~Lakatos,
\emph{Proofs and Refutations}, Cambridge University Press, 1976.

\bibitem{thurston}
W.~P.~Thurston,
On proof and progress in mathematics,
\emph{Bull.~Amer.~Math.~Soc.}~30 (1994), 161--177.

\bibitem{hamming}
R.~W.~Hamming,
The unreasonable effectiveness of mathematics,
\emph{Amer.~Math.~Monthly}~87 (1980), 81--90.

\bibitem{tao}
T.~Tao,
\emph{There's more to mathematics than rigour and proofs},
blog post, 2007.

\bibitem{mathlib}
The Mathlib Community,
\emph{Mathlib: A unified library of mathematics formalized in Lean},
2024.

\end{thebibliography}

\end{document}